Standard RLHF datasets (Anthropic-HH, WebGPT) cannot measure preference diversity — not because diversity was ignored, but because pairwise comparison formats are structurally incompatible with entropy measurement. Each example compares different response candidates (A1 vs B1, A2 vs B2, ...), preventing aggregation into preference distributions over fixed response sets. Entropy computation on Anthropic-HH yielded constant H = ln(2) $\approx$ 0.693 nats across all 100 sampled prompts (variance = 0) — not a measurement error, but a tautological artifact of the pairwise format.

\textbf{Bidirectional alignment} measures whether humans preserve critical evaluation capacity (preference diversity) when exposed to aligned AI, not just whether AI matches human preferences (agreement). We propose preference entropy as an information-theoretic proxy: Shannon entropy H = -$\Sigma$ $p_i$ log($p_i$) quantifies distribution diversity, with high entropy indicating diverse judgments (preserved critical evaluation) and low entropy indicating convergence (potential homogenization). This matters because high preference agreement (e.g., InstructGPT's 85\% win rate) could indicate successful alignment (appropriate for objective tasks like math, factual QA with clear ground truth) or problematic habituation on \textbf{subjective tasks} (creative writing, opinion questions, stylistic preferences where diverse preferences are legitimate) — existing metrics cannot distinguish these scenarios.

\textbf{Scope clarification:} Preference convergence is not universally problematic. For objective tasks with verifiable ground truth (e.g., arithmetic, factual questions), high agreement indicates successful learning rather than homogenization. Our concern applies primarily to subjective tasks where legitimate diversity should persist even after alignment. This distinction is essential to interpreting preference entropy as bidirectional alignment metric.

Our validation demonstrates that entropy computation is technically feasible (100\% success rate, correct mathematical range [0, ln(2)] for binary choices) but dataset format incompatibility prevents diversity measurement. We document a dataset structure trade-off: pairwise-unique formats (cost-efficient for reward model training) versus multi-annotator-fixed formats (enable entropy measurement at higher annotation cost). This explains why current RLHF benchmarks structurally cannot retrospectively measure diversity, even if desired.

We propose: (1) \textbf{Bidirectional alignment framework} introducing preference entropy as proxy for collective critical evaluation capacity (computable from existing preference data when structure supports it); (2) \textbf{Dataset structure documentation} identifying cost-measurability trade-off between pairwise-unique (Anthropic-HH: entropy-incompatible) and multi-annotator-fixed formats (OpenAI Summarization: entropy-measurable), tested on n=1 dataset with inference for others; (3) \textbf{Technical validation} demonstrating entropy computation feasibility (100\% success rate) but dataset format incompatibility (zero variance); (4) \textbf{Alternative measurement proxies} when multi-annotator data unavailable — response diversity entropy (model outputs), intra-annotator variance (longitudinal), entropy-regularized RLHF (algorithm modification).

Current RLHF benchmarks optimize annotation cost (pairwise comparisons) but sacrifice diversity measurability. Future benchmarks can address this via multi-annotator evaluation subsets (1K prompts $\times$ 20 annotators, \textasciitilde{}10-20\% overhead) or alternative proxies for retrospective analysis. The hypothesis mechanism (entropy collapse beyond performance-justified levels on subjective tasks) remains untested due to dataset limitation, not falsified.

\textbf{Keywords:} RLHF, preference diversity, bidirectional alignment, Shannon entropy, dataset structure, evaluation metrics
